\documentclass[12pt]{article}
\usepackage[utf-8]{inputenc}
\usepackage[spanish]{babel}
\usepackage{geometry}
\usepackage{graphicx}
\usepackage{hyperref}
\usepackage{listings}
\usepackage{xcolor}
\usepackage{fancyhdr}
\usepackage{amsmath}

\geometry{margin=1in}

% Configurar código
\lstset{
    language=go,
    basicstyle=\ttfamily\small,
    keywordstyle=\color{blue},
    commentstyle=\color{gray},
    stringstyle=\color{red},
    breaklines=true,
    tabsize=4,
    showstringspaces=false,
    numbers=left,
    numberstyle=\tiny\color{gray},
    backgroundcolor=\color{lightgray!20}
}

% Encabezado y pie de página
\pagestyle{fancy}
\fancyhf{}
\rhead{Analizador Sintáctico}
\lhead{\today}
\cfoot{\thepage}

\title{\textbf{Analizador Sintáctico} \\ \textit{Documentación y Explicación del Proyecto}}
\author{Analizador de Expresiones Matemáticas}
\date{\today}

\begin{document}

\maketitle

\newpage
\tableofcontents
\newpage

% ============================================
\section{Introducción}
% ============================================

Un \textbf{analizador sintáctico} (o \textit{parser}) es una herramienta fundamental en la teoría de compiladores. Su propósito es tomar una secuencia de símbolos (tokens) y validar si siguen las reglas de una gramática formal, construyendo una estructura de árbol sintáctico que representa la composición jerárquica del programa o expresión.

Este proyecto implementa \textbf{tres métodos diferentes} de análisis sintáctico para validar y analizar expresiones matemáticas simples con operadores de suma (+), multiplicación (*) y paréntesis.

\subsection{Objetivos del Proyecto}

\begin{itemize}
    \item Implementar tres algoritmos distintos de análisis sintáctico
    \item Proporcionar una interfaz gráfica interactiva y amigable
    \item Demostrar cómo funcionan diferentes estrategias de análisis
    \item Explicar paso a paso el proceso de análisis
\end{itemize}

% ============================================
\section{Conceptos Fundamentales}
% ============================================

\subsection{Gramática Formal}

La gramática utilizada en este proyecto es una \textbf{gramática libre de contexto (CFG)} que define las reglas para expresiones matemáticas válidas:

\begin{verbatim}
E  → T E'
E' → + T E' | ε
T  → F T'
T' → * F T' | ε
F  → ( E ) | num
\end{verbatim}

\textbf{Explicación:}
\begin{itemize}
    \item $E$ (Expresión): Representa una expresión completa
    \item $T$ (Término): Un término es un factor multiplicado por otros factores
    \item $F$ (Factor): Un factor es un número o una subexpresión entre paréntesis
    \item $E'$ y $T'$: Son producciones recursivas que permiten múltiples operaciones
    \item $\epsilon$ (Épsilon): Representa la producción vacía
\end{itemize}

\subsection{Tokens}

Los tokens son los símbolos básicos que el analizador reconoce:

\begin{itemize}
    \item \texttt{num}: Un número (cualquier secuencia de dígitos)
    \item \texttt{+}: Operador de suma
    \item \texttt{*}: Operador de multiplicación
    \item \texttt{(} y \texttt{)}: Paréntesis para agrupar
    \item \texttt{EOF}: Fin de entrada
\end{itemize}

% ============================================
\section{Métodos de Análisis Sintáctico}
% ============================================

Este proyecto implementa tres enfoques distintos: dos descendentes (top-down) y uno ascendente (bottom-up).

\subsection{1. Análisis Predictivo Recursivo}

\subsubsection{¿Qué es?}

Es un método \textbf{descendente} que implementa directamente cada regla de la gramática como una función recursiva. Comienza desde el símbolo inicial y desciende hacia los terminales.

\subsubsection{Cómo Funciona}

\begin{enumerate}
    \item \textbf{Lee un token} de la entrada
    \item \textbf{Decide qué regla aplicar} basándose en el token actual
    \item \textbf{Llama recursivamente} a las funciones correspondientes
    \item \textbf{Verifica} que cada token coincida con lo esperado
    \item \textbf{Retrocede} cuando encuentra una producción vacía ($\epsilon$)
\end{enumerate}

\subsubsection{Implementación en Go}

\begin{lstlisting}
// E → T E'
func (p *PredictivoRecursivo) E() bool {
    return p.T() && p.El()
}

// E' → + T E' | ε
func (p *PredictivoRecursivo) El() bool {
    if p.currentToken == "+" {
        return p.match("+") && p.T() && p.El()
    }
    return true // Producción vacía
}
\end{lstlisting}

\subsubsection{Ventajas}

\begin{itemize}
    \item Directo y fácil de entender
    \item Directamente relacionado con la gramática
    \item No requiere tabla predictiva precalculada
\end{itemize}

\subsubsection{Desventajas}

\begin{itemize}
    \item Limitado a gramáticas LL(1)
    \item Costo de llamadas recursivas
    \item Puede causar desbordamiento de pila con entradas grandes
\end{itemize}

% ============================================

\subsection{2. Análisis Predictivo No Recursivo}

\subsubsection{¿Qué es?}

Es un método \textbf{descendente} que utiliza una \textbf{tabla predictiva} precalculada y una \textbf{pila explícita} en lugar de recursión. Evita el costo de las llamadas de función.

\subsubsection{Cómo Funciona}

\begin{enumerate}
    \item \textbf{Inicializa una pila} con el símbolo inicial
    \item \textbf{Lee un token} de la entrada
    \item \textbf{Consulta la tabla predictiva} para el tope de la pila y el token actual
    \item \textbf{Ejecuta la acción}: desplazar, reducir o aceptar
    \item \textbf{Modifica la pila} según la acción
\end{enumerate}

\subsubsection{Tabla Predictiva}

La tabla predictiva $M[A, a]$ indica qué producción usar cuando:
\begin{itemize}
    \item $A$ es el símbolo no-terminal en el tope de la pila
    \item $a$ es el siguiente token de entrada
\end{itemize}

\subsubsection{Ventajas}

\begin{itemize}
    \item No usa recursión (más eficiente)
    \item Mayor control sobre el tamaño de la pila
    \item Separación clara entre datos y lógica
\end{itemize}

\subsubsection{Desventajas}

\begin{itemize}
    \item Requiere precalcular la tabla
    \item También limitado a gramáticas LL(1)
    \item Más complejo de debuggear
\end{itemize}

% ============================================

\subsection{3. Análisis Ascendente LR}

\subsubsection{¿Qué es?}

Es un método \textbf{ascendente} (bottom-up) que construye el árbol sintáctico desde las hojas hacia la raíz. Usa una tabla \textbf{LR(0)} que define un autómata de estados.

\subsubsection{Cómo Funciona}

\begin{enumerate}
    \item \textbf{Desplazamiento}: Agrega tokens a una pila hasta encontrar una "manija"
    \item \textbf{Reducción}: Cuando hay una "manija", aplica una regla de producción
    \item \textbf{GOTO}: Empuja el símbolo no-terminal resultante
    \item \textbf{Aceptación}: Cuando se reduce al símbolo inicial
\end{enumerate}

\subsubsection{Estados y Transiciones}

El analizador LR mantiene una máquina de estados donde:
\begin{itemize}
    \item \textbf{Estados}: Representan configuraciones de análisis
    \item \textbf{Transiciones}: Definidas por la tabla LR
    \item \textbf{Acciones}: Shift (desplazar), Reduce (reducir), Accept (aceptar)
\end{itemize}

\subsubsection{Ventajas}

\begin{itemize}
    \item Maneja gramáticas más complejas (no solo LL(1))
    \item Detecta errores más rápidamente que métodos descendentes
    \item Puede procesar más lenguajes
\end{itemize}

\subsubsection{Desventajas}

\begin{itemize}
    \item Mucho más complejo de entender
    \item La construcción de la tabla LR es tedioso
    \item Mayor consumo de memoria para la tabla
\end{itemize}

% ============================================
\section{Estructura del Código}
% ============================================

\subsection{Archivos Principales}

\subsubsection{main.go}

Contiene el servidor web HTTP y la interfaz gráfica:

\begin{itemize}
    \item \texttt{handleIndex()}: Sirve la página HTML principal
    \item \texttt{handleAnalizar()}: Punto de entrada de API REST
    \item \texttt{analizarPredictivoRecursivo()}: Llama al método y captura resultados
    \item \texttt{analizarPredictivoNoRecursivo()}: Análisis no-recursivo
    \item \texttt{analizarLR()}: Análisis ascendente LR
    \item \textbf{htmlContent}: Contiene la interfaz gráfica completa en HTML/CSS/JavaScript
\end{itemize}

\subsubsection{predictivo\_recursivo.go}

Implementación del método predictivo recursivo:

\begin{itemize}
    \item \texttt{PredictivoRecursivo}: Estructura con estado del analizador
    \item \texttt{E()}, \texttt{El()}, \texttt{T()}, \texttt{Tl()}, \texttt{F()}: Funciones para cada regla gramatical
    \item \texttt{nextToken()}: Tokenización de entrada
    \item \texttt{Parse()}: Función principal de análisis
    \item \texttt{pasos}: Lista de pasos ejecutados para explicación
\end{itemize}

\subsubsection{desendente\_no\_recursivo.go}

Implementación del método predictivo no-recursivo:

\begin{itemize}
    \item \texttt{PredictivoNoRecursivo}: Con tablas predictivas
    \item \texttt{Pila}: Estructura para mantener símbolos
    \item \texttt{tablaAccion}: Tabla de decisiones sintácticas
    \item \texttt{Parse()}: Análisis iterativo con pila
\end{itemize}

\subsubsection{ascendente\_lr.go}

Implementación del método ascendente LR:

\begin{itemize}
    \item \texttt{AnalizadorLR}: Con tabla de acción y GOTO
    \item \texttt{EstadoLR}: Estados del autómata
    \item \texttt{ArbolSintactico}: Estructura del árbol construido
    \item \texttt{construirTablas()}: Genera la tabla LR(0)
    \item \texttt{Parse()}: Algoritmo shift-reduce
\end{itemize}

\subsection{Flujo de Ejecución}

\begin{verbatim}
Usuario abre navegador
        ↓
Interfaz web (HTML/CSS/JavaScript)
        ↓
Usuario introduce expresión y selecciona método
        ↓
Petición POST a /api/analizar
        ↓
handleAnalizar() recibe datos
        ↓
    ↙              ↓              ↘
Recursivo    No Recursivo        LR
    ↓              ↓              ↓
Crea objeto  Crea objeto      Crea objeto
    ↓              ↓              ↓
Llama Parse()  Llama Parse()   Llama Parse()
    ↓              ↓              ↓
Retorna        Retorna          Retorna
resultados     resultados       resultados
    ↓              ↓              ↓
    ↖              ↓              ↗
        Respuesta JSON
             ↓
    JavaScript procesa
             ↓
    Muestra resultados
\end{verbatim}

% ============================================
\section{Interfaz Gráfica}
% ============================================

\subsection{Descripción General}

La interfaz web es responsiva y moderna, construida con:

\begin{itemize}
    \item \textbf{HTML5}: Estructura de la página
    \item \textbf{CSS3}: Estilos con gradientes y animaciones
    \item \textbf{Bootstrap Icons}: Iconografía visual
    \item \textbf{JavaScript vanilla}: Interactividad sin dependencias
\end{itemize}

\subsection{Componentes Principales}

\subsubsection{Panel de Configuración}

\begin{itemize}
    \item Campo de entrada para la expresión
    \item Botones de ejemplo rápido
    \item Selector de método de análisis
    \item Botones de análisis y limpieza
\end{itemize}

\subsubsection{Panel de Información}

\begin{itemize}
    \item Descripción de cada método
    \item Características diferenciales
    \item Información educativa
\end{itemize}

\subsubsection{Panel de Resultados}

\begin{itemize}
    \item Mensaje de éxito o error
    \item Tabla de pasos del análisis
    \item Árbol sintáctico resultante
    \item Explicaciones detalladas del proceso
\end{itemize}

\subsection{Paleta de Colores}

\begin{itemize}
    \item \textbf{\#2D261D} (Marrón muy oscuro): Fondo principal
    \item \textbf{\#A79675} (Beige/Marrón): Acentos y gradientes
    \item \textbf{\#C1B18B} (Dorado): Detalles y botones
    \item \textbf{\#DCCFAA} (Crema): Texto en encabezados
    \item \textbf{\#F5F3EF} (Crema muy clara): Cards
\end{itemize}

% ============================================
\section{Ejemplos de Uso}
% ============================================

\subsection{Ejemplo 1: Expresión Válida Simple}

\textbf{Entrada:} \texttt{3+5*2}

\textbf{Procesamiento:}

\begin{enumerate}
    \item Tokenización: [3, +, 5, *, 2, EOF]
    \item Método Recursivo: Aplica reglas E → T E' → ... hasta consumir todos los tokens
    \item Método No Recursivo: Consulta tabla y reduce según reglas
    \item Método LR: Desplaza tokens, reduce producs, hasta llegar a S (éxito)
\end{enumerate}

\textbf{Resultado:} ✅ Análisis exitoso

\subsection{Ejemplo 2: Expresión con Paréntesis}

\textbf{Entrada:} \texttt{(3+5)*2}

\textbf{Procesamiento:}

Los paréntesis fuerzan una subexpresión que es procesada recursivamente, permitiendo que el analizador comprenda la precedencia correcta de operadores.

\textbf{Resultado:} ✅ Análisis exitoso

\subsection{Ejemplo 3: Expresión Inválida}

\textbf{Entrada:} \texttt{3++5}

\textbf{Procesamiento:}

Cuando se encuentra el segundo \texttt{+}, el analizador espera un número o paréntesis (un factor), no otro operador.

\textbf{Resultado:} ❌ Error de sintaxis con explicación

% ============================================
\section{Cómo Ejecutar el Programa}
% ============================================

\subsection{Requisitos}

\begin{itemize}
    \item Go 1.19 o superior
    \item Un navegador web moderno
\end{itemize}

\subsection{Pasos de Ejecución}

\begin{enumerate}
    \item Clonar o descargar el proyecto
    \item Abrir una terminal en el directorio del proyecto
    \item Compilar: \texttt{go build -o app.exe main.go *.go}
    \item Ejecutar: \texttt{./app.exe}
    \item Abirir navegador: \texttt{http://localhost:3000}
\end{enumerate}

% ============================================
\section{Comparativa de Métodos}
% ============================================

\begin{table}[h]
\centering
\begin{tabular}{|l|c|c|c|}
\hline
\textbf{Característica} & \textbf{Recursivo} & \textbf{No Recursivo} & \textbf{LR} \\
\hline
Dirección & Descendente & Descendente & Ascendente \\
\hline
Recursión & Sí & No & No \\
\hline
Tabla Predictiva & No & Sí & Sí (LR) \\
\hline
Gramáticas LL(1) & Sí & Sí & No (más complejas) \\
\hline
Complejidad & Baja & Media & Alta \\
\hline
Velocidad & Media & Rápido & Rápido \\
\hline
Consumo Memoria & Alto (stack) & Medio & Medio \\
\hline
Fácil de Entender & Muy & Medio & No \\
\hline
Producción & No & Sí & Sí \\
\hline
\end{tabular}
\caption{Comparativa entre métodos de análisis sintáctico}
\end{table}

% ============================================
\section{Conclusiones}
% ============================================

Este proyecto demuestra de forma práctica cómo funcionan los tres enfoques principales de análisis sintáctico:

\begin{enumerate}
    \item \textbf{Predictivo Recursivo}: Ideal para educación y comprensión fundamental
    \item \textbf{Predictivo No Recursivo}: Balance entre simplicidad y eficiencia
    \item \textbf{Ascendente LR}: Estándar industrial para compiladores reales
\end{enumerate}

Cada método tiene sus ventajas y desventajas, y la elección depende de:
\begin{itemize}
    \item La complejidad de la gramática
    \item Los requisitos de rendimiento
    \item El mantenimiento del código
    \item La portabilidad
\end{itemize}

La interfaz gráfica proporciona una forma interactiva y educativa de entender cómo funcionan estos algoritmos sin necesidad de estudiar código complejo.

% ============================================
\section{Referencias}
% ============================================

\begin{enumerate}
    \item Aho, A. V., Lam, M. S., Sethi, R., \& Ullman, J. D. (2006). 
    \textit{Compilers: Principles, Techniques, and Tools} (2nd ed.). Pearson Education.
    
    \item Appel, A. W., \& Palsberg, J. (2002). 
    \textit{Modern Compiler Implementation in Java} (2nd ed.). Cambridge University Press.
    
    \item Language Specification Go: \texttt{https://golang.org/ref/spec}
    
    \item Bootstrap Icons: \texttt{https://icons.getbootstrap.com/}
\end{enumerate}

% ============================================
\section{Apéndice: Notación}
% ============================================

\subsection{Símbolos Gramaticales}

\begin{itemize}
    \item Letras mayúsculas: Símbolos no-terminales (E, T, F)
    \item Letras minúsculas: Terminales (números, operadores)
    \item $\epsilon$: Producción vacía
    \item $|$: Alternativa (o)
    \item $\rightarrow$: Produce
\end{itemize}

\subsection{Notación de Máquinas de Estado}

\begin{itemize}
    \item Estados: Números (0, 1, 2, ...)
    \item Transiciones: Etiquetadas con símbolos
    \item Estados de aceptación: Indicados con \textbf{acc}
    \item Estados de error: Sin transición disponible
\end{itemize}

\end{document}
